\begin{claim}{D.2.16}{Proved}{A fourth-stage obstruction to the Hopf coefficient}
Let $B=K(\FF_2,2)$, $X=\Sigma B$, $Y=\Sigma^2B$, and
$C_mX=\Omega^m\Sigma^mX$. Write $a\in H_2(B;\FF_2)$ for
the bottom class, $c\in H_4(B;\FF_2)$ for the class dual to
the square of the fundamental cohomology class, and $C=\sigma c$.
For every ordinary pointed map
\[
 G:C_4X\longrightarrow T=\Omega^2Y,
\]
the coefficient of the top $E_2$ generator $\xi_1a$ in
$G_*(C)$ is zero. In particular, the map
\[
 h_0=\eta_2\wedge\id_B:\Sigma^3B\longrightarrow\Sigma^2B
 \tag{D.2.16a}
\]
does not extend across
$\Sigma^2X\to\Sigma^2C_4X$, for any choice of ordinary pointed
extension. Consequently it does not extend across
$\Sigma^2X\to\Sigma^2QX$.

This excludes the particular homology-zero Hopf attachment whose
cofiber is $\mathbb CP^2\wedge B$. It does not assert survival
through all subsequent $Q\Sigma B$-nullification attachments.
\end{claim}
\begin{proof}
All homology and cohomology in this proof have coefficients $\FF_2$.
Put $T=\Omega^2Y$, and let $g_0:X\to T$ be the double adjunct of
$h_0$. We use the homology labels $a,b,c,d,e,f,x$ of D.2.13,
dual respectively to $u,v,u^2,w,uv,u^3,v^2$ in
$H^*B$, where
\[
 |u|=2,\qquad v=\operatorname{Sq}^1u,\qquad
 w=\operatorname{Sq}^2v,\qquad t=\operatorname{Sq}^4w.
 \tag{D.2.16b}
\]
Their degrees are $2,3,5,9$ respectively. Serre's calculation
makes these independent polynomial generators; in particular
$t\ne0$ \cite[\S2, nos.~8--9, pp.~203--204]{I_Serre1953}.
Write $\beta=\operatorname{Sq}_*^1$ on homology,
$\xi=\xi_1$ for the top $E_2$ homology operation, and
$\{\ ,\ \}=\lambda_1$ for the Browder bracket.

\medskip\noindent
\textbf{The forced failure of multiplicativity.}
Set $C=\sigma c\in H_5X$, and use the same letter for its image
in the loop stages. The calculation of D.2.13 gives
\[
 g_{0*}(C)=\xi a.
 \tag{D.2.16c}
\]
Let $L\in H^{10}T$ be the polynomial-basis functional dual to
$(\xi a)^2$. The fourth-stage calculation in D.2.13 proves, for
\emph{every} ordinary map $G:C_4X\to T$,
\[
 \langle L,G_*(C^2)\rangle=0.
 \tag{D.2.16d}
\]
Here $C^2$ is the Pontryagin square. Explicitly,
$\beta\operatorname{Sq}_*^2\xi_3C=C^2$ in the fourth stage,
where the top Nishida relation's Browder term has zero
Bockstein. On the target, the coefficient of $(\xi a)^2$ in
$\beta\operatorname{Sq}_*^2H_{13}T$ is zero. Thus
\textup{(D.2.16d)} uses only Steenrod naturality for $G$.

\medskip\noindent
\textbf{Exact cohomology identities in the target.}
In addition to $L$, define polynomial-basis dual classes
\[
\begin{array}{c|ccccc}
\text{class}&U&V&W&L&L_b\\
\text{degree}&5&7&9&10&14\\
\text{dual monomial}&\xi a&\xi b&\xi c&(\xi a)^2&(\xi b)^2.
\end{array}
 \tag{D.2.16e}
\]
We will use the following four identities:
\[
 \operatorname{Sq}^2L=UV,\qquad
 \operatorname{Sq}^4L=L_b+UW,\qquad
 \operatorname{Sq}^2L_b=0,\qquad
 \operatorname{Sq}^2\operatorname{Sq}^1L=0.
 \tag{D.2.16f}
\]
The last is the target coefficient identity of D.2.13. We check
all possibilities for the other three.

The free $E_2$ homology theorem describes $H_*T$ as a polynomial
algebra on a free restricted degree-one Lie basis on the labels
of $\widetilde H_*B$; see
\cite[Chapter III, Theorems~3.1--3.2 and Lemma~3.8]{I_CohenLadaMay1976}.
Every pure top generator $\xi z$ has odd degree. Steenrod
operations send labels to labels, and send bracket expressions
to the bracket ideal. They cannot send a positive-degree
factor to a nonzero scalar. Thus a pure two-factor polynomial
coefficient in an even-degree Steenrod image must come from an
initial product of two Lie generators. No top generator exists
in that even input degree.

The exact top Nishida formula gives
\[
 \operatorname{Sq}_*^2\xi b=\xi a,\qquad
 \operatorname{Sq}_*^2\xi c=\{c,a\},\qquad
 \operatorname{Sq}_*^4\xi c=\xi a.
 \tag{D.2.16g}
\]
For the last identity, use
$\beta c=0$, $\operatorname{Sq}_*^2c=a$, and
$\operatorname{Sq}_*^3c=\operatorname{Sq}_*^4c=0$.
The surviving ordinary term is $Q^3a=\xi a$ and all the
Browder terms vanish. These are the original formulas of
\cite[Chapter III, Theorems~1.2--1.4, pp.~216--219,
especially pp.~217--218]{I_CohenLadaMay1976}.

There is no degree-six Bockstein preimage of the indecomposable
$\xi a$. The full degree-six basis is
\[
 f,x,ac,b^2,a^3,\{a,b\},
\]
with Bockstein images $0,d+e,0,0,0,0$. There is also no
degree-eight Bockstein preimage of the indecomposable $\xi b$:
even-degree indecomposables are labels or brackets, and their
Bocksteins cannot become a pure top generator; a decomposable
class has decomposable Bockstein.

The coefficient of $(\xi a)^2$ in
$\operatorname{Sq}_*^2H_{12}T$ is consequently exactly the
coefficient of $\xi a\cdot\xi b$ in its input. The partition
$1+1$ would require the absent degree-six Bockstein preimage.
Likewise the coefficient of $(\xi a)^2$ in
$\operatorname{Sq}_*^4H_{14}T$ is the sum of the input
coefficients of $(\xi b)^2$ and $\xi a\cdot\xi c$.
The partitions $1+3$ and $3+1$ again require the absent
Bockstein preimage. The partitions $2+2$ and $0+4$ are
precisely those displayed in \textup{(D.2.16g)}.

These two product-coefficient functionals are exactly the cup
products $UV$ and $UW$. This point requires the coalgebra, since
$c$ is not primitive. Coproducts of labels contain only labels;
coproducts of Browder generators retain a bracket in at least
one factor; and there are no even-degree top generators in
degrees twelve or fourteen. For a product to have the two
specified indecomposable top generators in its tensor factors,
its factors must therefore be exactly the indicated two pure
top generators. A product of lower label components cannot
equal an indecomposable $\xi a$, $\xi b$, or $\xi c$.
This proves the first two identities of \textup{(D.2.16f)}
as exact cohomology identities, rather than just as projections
on a polynomial subspace.

Finally, a degree-sixteen class could contribute $(\xi b)^2$
after $\operatorname{Sq}_*^2$ only through a degree-nine
square-operation preimage of $\xi b$ or a degree-eight
Bockstein preimage. The latter was excluded above. The only
pure top generator in degree nine is $\xi c$, whose image in
\textup{(D.2.16g)} is a bracket. Products and the other
indecomposables cannot supply the required preimage. This gives
$\operatorname{Sq}^2L_b=0$ and completes the check of
\textup{(D.2.16f)}.

\medskip\noindent
\textbf{An actual map representing the multiplication deviation.}
Suppose that $G:C_4X\to T$ has coefficient one on $\xi a$
in $y=G_*(C)$. Set $g=Gj_4$, where
$j_4:X\to C_4X$ is the unit. On $X\times X$, form the two maps
\[
 F=G\mu_4(j_4\times j_4),\qquad
 H=\mu_T(g\times g),
 \tag{D.2.16h}
\]
where $\mu_4$ and $\mu_T$ are the loop multiplications.
They agree up to a pointed homotopy on $X\vee X$.
Their pointwise difference, formed by multiplying $F$ with the
pointwise inverse of $H$ in the H-group $T$, is therefore null
on that wedge. Choosing the nullhomotopy and using the pointed
cofiber universal property produces an ordinary map
\[
 \delta:X\wedge X\simeq\Sigma^2(B\wedge B)\longrightarrow T.
 \tag{D.2.16i}
\]
One may use a loop-group model for strict inverses. No
multiplicativity of $G$, coherent symmetry of $\delta$, or
factorization through a homotopy-orbit square is assumed.

The classes $C$ and $y$ are primitive. On $C\otimes C$, the homology of
the pointwise difference is computed by convolution. Its two
cross terms are both $y^2$ and cancel in characteristic two.
The antipode fixes the primitive $y$ and its square. In the
polynomial algebra $H_*T$, the coefficient of $(\xi a)^2$
in $y^2$ is exactly the coefficient of $\xi a$ in $y$, namely one.
Equation \textup{(D.2.16d)} therefore gives
\[
 \langle L,\delta_*(C\otimes C)\rangle
  =\langle L,G_*(C^2)\rangle+\langle L,y^2\rangle=1.
 \tag{D.2.16j}
\]
The quotient $X\times X\to X\wedge X$ identifies the K\"unneth
cross summand with the reduced homology of the smash product.
Thus the chosen nullhomotopy on the wedge does not alter this
cross-homology calculation.

\medskip\noindent
\textbf{The finite contradiction on the smash source.}
Let $z\in\widetilde H^8(B\wedge B)$ be the double desuspension
of $\delta^*L$. By \textup{(D.2.16j)}, its coefficient of
$u^2\otimes u^2$ is one. All positive cup products on
$X\wedge X$ vanish, since it is a double suspension.
Pulling back \textup{(D.2.16f)} and using compatibility of
Steenrod operations with suspension yields
\[
 \operatorname{Sq}^2z=0,\qquad
 \operatorname{Sq}^2\operatorname{Sq}^1z=0,\qquad
 \operatorname{Sq}^2\operatorname{Sq}^4z=0.
 \tag{D.2.16k}
\]
For the third equality, first pull back
$\operatorname{Sq}^4L=L_b+UW$, and then use
$\operatorname{Sq}^2L_b=0$.

The complete degree-eight basis of
$\widetilde H^*(B\wedge B)$ is
\[
\begin{array}{lll}
 u\otimes u^3,&u\otimes v^2,&v\otimes w,\\
 v\otimes uv,&u^2\otimes u^2,&w\otimes v,\\
 uv\otimes v,&u^3\otimes u,&v^2\otimes u.
\end{array}
 \tag{D.2.16l}
\]
Cartan gives
\[
 \ker(\operatorname{Sq}^2\text{ in degree }8)
       =\langle z_0,z_1\rangle,
 \quad
 \begin{cases}
 z_0=u^2\otimes u^2+u\otimes v^2+v^2\otimes u,\\
 z_1=v\otimes w+w\otimes v.
 \end{cases}
 \tag{D.2.16m}
\]
Here is an exhaustive coefficient check. The terms
$u\otimes u^4$ and $u^4\otimes u$ force the coefficients of
$u\otimes u^3$ and $u^3\otimes u$ to vanish. The unique terms
$w\otimes uv$ and $uv\otimes w$ force the coefficients of
$v\otimes uv$ and $uv\otimes v$ to vanish. The remaining
images force the three coefficients in $z_0$ to agree and the
two in $z_1$ to agree. This proves \textup{(D.2.16m)}.

Furthermore,
\[
 \operatorname{Sq}^1z_0=\operatorname{Sq}^1z_1
       =v\otimes v^2+v^2\otimes v,
\]
\[
 \operatorname{Sq}^2\operatorname{Sq}^1z_0
 =\operatorname{Sq}^2\operatorname{Sq}^1z_1
       =w\otimes v^2+v^2\otimes w\ne0.
 \tag{D.2.16n}
\]
The first two conditions in \textup{(D.2.16k)}, with the required
coefficient one, therefore force
\[
 z=z_0+z_1,
 \qquad\operatorname{Sq}^1z=0.
 \tag{D.2.16o}
\]
The Adem relation
$\operatorname{Sq}^2\operatorname{Sq}^4
 =\operatorname{Sq}^6+\operatorname{Sq}^5\operatorname{Sq}^1$
now gives
\[
 \operatorname{Sq}^2\operatorname{Sq}^4z
       =\operatorname{Sq}^6z\ne0.
 \tag{D.2.16p}
\]
Indeed the coefficient of $w\otimes t$ is one: it comes from
$\operatorname{Sq}^2v\otimes\operatorname{Sq}^4w$ on the term
$v\otimes w$, and none of the other four terms of
\textup{(D.2.16o)} can contribute that basis element.
This contradicts the third condition in \textup{(D.2.16k)}.

Thus no $G$ can have the asserted nonzero $\xi a$ coefficient.
Equation \textup{(D.2.16c)} and double adjunction prove the
claimed fourth-stage nonextension of $h_0$. Restricting a
hypothetical full $QX$-extension to this stage gives the same
contradiction.
The finite coefficient calculation was independently verified
by exhaustive enumeration of the nine-dimensional source space;
the displayed proof does not depend on that computation.
\end{proof}
